\section{Conclusions}
\label{sec:conclusion}

This study presented an experimental investigation of paraffin-based PCM (phase change temperature 42\textdegree C) infiltrated into AlSiMg Gyroid TPMS lattice structures fabricated by laser powder bed fusion. The thermal response was systematically characterized under three heating power levels (10~W, 20~W, 30~W) using a nine-sensor array organized into three spatial groups. The key findings are:

\begin{enumerate}[nosep]
  \item The Gyroid TPMS lattice provides effective thermal conductivity enhancement, with the temperature spread between the innermost (Group~A) and outermost (Group~C) sensor groups limited to 17.1\textdegree C at 10~W---substantially smaller than typical bare PCM systems.

  \item A consistent radial thermal stratification ($T_1 > T_{2\text{-}5} > T_{6\text{-}9}$) was observed across all power levels, reflecting the heat flow direction from the central heater through the TPMS lattice/PCM composite.

  \item The heating power significantly affects the phase change dynamics: at 10~W, a clear phase change plateau near 42\textdegree C was observed for 5--8 minutes; at 20~W, the plateau compressed to 3--4 minutes; at 30~W, the plateau was barely discernible due to the high heating rate.

  \item The experiment duration decreased from 25.4~min (10~W) to 11.1~min (20~W) to 8.5~min (30~W), with the outer sensor groups (B and C) showing more monotonic temperature rise trends than the innermost group (A).

  \item Systematic outlier analysis revealed that $T_6$ was consistently identified as an outlier in Group~C, while the Group~B outlier varied with heating power, suggesting power-dependent convection patterns in the molten PCM.
\end{enumerate}

The present work establishes the Gyroid TPMS/PCM baseline for an ongoing comparative study. Future experiments will extend the test matrix to Primitive and IWP topologies at all three heating powers (Table~\ref{tab:test_matrix}), enabling direct comparison of the thermal performance across TPMS families. The comparative data will inform the selection of optimal TPMS topologies for specific thermal management applications, considering trade-offs between thermal response rate, temperature uniformity, and energy storage capacity.

Several limitations of the current study should be noted. First, only the heating (melting) phase has been characterized; the cooling (solidification) behavior of the TPMS/PCM composite remains to be investigated. Second, the current analysis is based solely on temperature measurements; future work should incorporate heat flux measurements and energy balance calculations to quantify the effective thermal conductivity enhancement factor. Third, the effect of natural convection within the molten PCM channels of the TPMS lattice has not been isolated; complementary numerical simulations using the enthalpy-porosity method~\cite{ebrahimi2019sensitivity} are planned to decompose the conduction and convection contributions.
